\scanpage{37}
Montrons que l'axiome $D=d)+d')$ implique

\uuline{Corollaire} 1 (moyennant l'axiome D) pour \unsure{tte|tt} $s\in S$,
\deleted{\unsure{on}} $p(s)=\mathrm{Card}\,(A\{s\}=\{a\in A\mid a\geq s\}$ \deleted{$\geq3$} $\unsure{\geq3|\geq2}$ et pour
tout $f\in F$, $q(f)=\mathrm{Card}\,(A\unsure{\}|\{}f\}=\{a\in A\mid a\leq f)\geq3$.

En effet, \deleted{si on avait} $p(s)\ne0$, en vertu de
c). Si on avait $p(s)=1$, on aurait par c) une
seule face incidente à l'\inserted{unique} arête $a$ incidente à $s$,
% slanted
\margin{\uline{NB} $p(s)\geq2$ est \unsure{indépendant|indépendt} de D, de \unsure{$=$|m\textsuperscript{me}} $q(f)\geq2$.}
contrairement à b). Si on avait $p(s)=2$, en vertu
de c) les deux arêtes issues de $s$ seraient incidentes
simultanément aux deux faces incidentes à $s$,
contrairement à $d'_1)$. Donc $p(s)\geq3$. La relation
$q(f)\geq3$ s'établit de même.

\uline{Corollaire 2} (moyennant l'axiome D) : toute face
$f\in F$ est le sup de l'ens de ses sommets,
et le sup de l'ens de ses arêtes.

Tout sommet est le inf de l'ens des
faces qui y sont incidentes, et le
Inf de l'ens des arêtes incidentes.
% archive sheet number: 36
